\section{Related Work}

\paragraph{Text detoxification.}
Text detoxification has been studied as controlled rewriting, style transfer, and
safety-oriented generation. Prior work commonly evaluates content preservation,
reference similarity, fluency, and toxicity reduction. Our setting differs by focusing
on multilingual paired toxic-to-detox rewrites and by separating preservation proxies
from independent non-toxicity evidence. We will cite the corpus source and the most
relevant multilingual detoxification baselines in the final literature pass.

\paragraph{Mechanistic interpretability and probes.}
Linear probes are widely used to test whether neural representations encode a target
property. However, probe separability does not by itself show that a representation is
causally used by a model. Our analysis therefore distinguishes pooled decodability,
cross-language transfer, and intervention effects. We will add citations to
representation-probing methodology, causal-probing critiques, and mechanistic
interpretability work in the final version.

\paragraph{Sparse autoencoders and activation steering.}
SAEs provide overcomplete sparse dictionaries for transformer activations, while
activation steering modifies an internal direction during generation. These methods
have motivated work on semantic features, behavior control, and safety. We adapt them
to paired multilingual toxic-to-detox transformations and explicitly report whether an
artifact supports descriptive feature ranking, exploratory steering behavior, or a
controlled causal claim.

\paragraph{Multilingual transfer and low-resource evaluation.}
Cross-lingual transfer often degrades for low-resource languages because of data
coverage, tokenization, morphology, and linguistic variation. We evaluate a specific
higher-resource-to-lower-resource condition: probes are trained on Amharic, Arabic,
German, English, Spanish, Hindi, and Ukrainian and tested on isiXhosa and Yoruba.
This is not a generic claim about all cross-lingual transfer.
